\documentclass[12pt,a4paper]{extarticle}
\usepackage[margin=2cm]{geometry}

\def\code#1{\texttt{#1}}

\begin{document}

\paragraph{Theory 1}

\subparagraph{Q-learning update} It is done at every timestep and consists in \[Q(s,a) \leftarrow Q(s,a) + \alpha(R + \gamma max_aQ(s',a) - Q(s,a))\] so, knowing that $(s,a,r,s') = (1,1,5,2), \alpha = 0.1$ and $\lambda = 0.5$ it holds that \[Q(s,a) = Q(1,1) + 0.1(5 + 0.5max_aQ(2,a)-Q(1,1)) = \] \[1 + 0.1(5+ 0.5*4 - 1) = 1.6\]

\subparagraph{Sarsa update} It is done at every timestep and consists in \[Q(s,a) \leftarrow Q(s,a) + \alpha(R + \gamma Q(s',a') - Q(s,a))\] so, knowing that $(s,a,r,s') = (1,1,5,2), a'=1$, it holds that \[Q(s,a) = Q(1,1) + 0.1(5 + 0.5*Q(2,1)-Q(1,1)) = \] \[1 + 0.1(5+ 0.5*3 - 1) = 1.55\]

\paragraph{Theory 2}

To compute the number of updates in an episode, it is convenient to first write the lambda return in a recursive form. We know that \[G_t^{\lambda} = (1-\lambda)\sum_{n=1}^{\infty} \lambda^{n-1} G_{t:t+n}\] which means that \[ G_t^{\lambda} = (1-\lambda)[G_{t:t+1}+\lambda G_{t:t+2} + \lambda^2 G_{t:t+3} + ...]\] by expanding the n-step return, we obtain
\[ G_t^{\lambda} = (1-\lambda)[(R_{t+1} + \gamma G_{t+1:t+1})+ \lambda (R_{t+1} + \gamma G_{t+1:t+2}) + \lambda^2 (R_{t+1} + \gamma G_{t+1:t+3}) + ...]\] now isolating $R_{t+1}$ we obtain \[ G_t^{\lambda} = (1-\lambda)[(R_{t+1} + \lambda R_{t+1} + \lambda^2 R_{t+1} + ...) + \gamma (G_{t+1:t+1}+ \lambda G_{t+1:t+2} + \lambda^2 G_{t+1:t+3} + ...)]\] the sum of rewards in the first parenthesis is a geometric series, so \[ G_t^{\lambda} = (1-\lambda)[\frac{R_{t+1}}{1-\lambda} + \gamma (G_{t+1:t+1}+ \lambda G_{t+1:t+2} + \lambda^2 G_{t+1:t+3} + ...)]\] taking out the fraction \[ G_t^{\lambda} = R_{t+1} + (1-\lambda)[\gamma (G_{t+1:t+1}+ \lambda G_{t+1:t+2} + \lambda^2 G_{t+1:t+3} + ...)] = \]\[R_{t+1} + (1-\lambda)\gamma (G_{t+1:t+1}) + (1-\lambda)\gamma[\lambda G_{t+1:t+2} + \lambda^2 G_{t+1:t+3} + ...)]\] by isolating $\lambda$, the third term of this sum yields $\gamma \lambda G_{t+1}^\lambda$, so finally \[G_{t}^\lambda = R_{t+1} + (1-\lambda)\gamma (G_{t+1:t+1}) +  \gamma \lambda G_{t+1}^\lambda\] where $G_{t+1:t+1} = \hat{v}(S_{t+1},\omega_{t})$, so basically it represents an approximated value function. Once the lambda return is written in recursive form, we can substitute this expression in the "montecarlo error" which is accumulated in an entire episode of temporal difference updates. \[G_t^\lambda - \hat{v}(S_t) = R_{t+1} + \gamma(1-\lambda)\hat{v}(S_{t+1}) +  \gamma \lambda G_{t+1}^\lambda - \hat{v}(S_t)\] by summing and subtracting $\gamma \hat{v}(S_{t+1})$ we obtain \[G_t^\lambda - \hat{v}(S_t) = \delta_t + \gamma \lambda G_{t+1}^\lambda + \gamma(1-\lambda)\hat{v}(S_{t+1}) - \gamma \hat{v}(S_{t+1})\] where $\delta_t = R_{t+1} + \gamma \hat{v}(S_{t+1}) - \hat{v}(S_t)$ is a one-step temporal difference error. Now, by summing and subtracting $\gamma\lambda\hat{v}(S_{t+1})$ we obtain \[G_t^\lambda - \hat{v}(S_t) = \delta_t + \gamma \lambda (G_{t+1}^\lambda - \hat{v}(S_{t+1}))\] where, exploiting the recursion we obtain \[G_{t}^\lambda - \hat{v}(S_{t}) = \delta_t + \gamma  \lambda[\delta_{t+1} + \gamma \lambda (G_{t+2}^\lambda - \hat{v}(S_{t+2}))]\] and finally by induction this yields that \[G_{t}^\lambda - \hat{v}(S_{t}) = \sum_{k=t}^\infty (\gamma \lambda)^{k-t} \delta_k\] if the episode terminates in T, this sum will be truncated, and so the single updates accumulated in the episode are of number ($\alpha$ is a constant) \[\sum_{k=t}^{T-1} (\gamma \lambda)^{k-t} \delta_k\] which is exactly what we wanted to demonstrated. So we showed that the entire error between the current value function estimate and the full lambda return can be expressed as a sum of temporal difference errors, weighted by the eligibility traces.


\paragraph{Practice 1}
The problem is to implement the Sarsa($\lambda$) for the frozen lake environment. So this is basically a multi-step learning algorithm and the number of steps considered (so the backward states) is related to lambda. Practically, for every episode the Q-table and the eligibility traces. The first is set to zero (or random) at the beginning of he learning process, while the latter is put to zero before every episode. What we do at every step is take the temporal difference error and give it to the q-function, in order to update it, weighing this error with a learning rate and the eligibility traces. But before updating the Q function, at every step we accumulate the visit at the current state in the eligibility traces, which means that the q functions contains values proportional to the number of visits and the temporal difference error for every state action pair. What is noticeable in the simulation is that the elf learns to step away from the pits by turning the back to them as soon as he is in their proximity. That's due to the fact that, since the terrain is slippery, going to a direction might cause with a certain probability to go to an orthogonal direction wrt the intentional one. Thus, if the pit is at the back of the elf, he surely won't step into it. Naturally, if there are no pits in sight, the Q function imposes to reach the goal as quickly as possible.

\paragraph{Practice 2}
The problem here is to implement the TD($\lambda$) Q-Learning algorithm with linear value function approximation in the mountain car environment. This problem can subdivided into two subproblems: how to represent the value function (in this specific case it will be a value-action function) and how to make the agent learn.
\subparagraph{First problem: RBF} For the approximated q-function representation, it is asked to encode the state with a radial basis function, which consists in characterizing each feature as a gaussian function which takes in input the distance between the state and a feature center, and the width of the feature. For this purpose, we need to first construct the features' center states, which in this case (the state space is of dimension 2) are basically some equally distributed points in 2D. The borders of the region considered ( [0,1]x[0,1]) don't contain any centers because, by normalizing the state, there won't be any out the region considered. The \code{encode} function will return an encoded state (a column vector) with as many features as centers taken for the rbf. TO generalize this behavior, one could implement multiple tilings. I implemented it, but didn't notice a bump in performance.
\subparagraph{Second problem: Q-Learning with eligibility traces} After having encoded the state, the agent needs to learn at every step. Intuitively, in the update of the weights we just assign the temporal difference  error (which in Q-learning takes the max of the Q of $s'$) backwards to the previous states by multiplying it with the eligibility traces. So the traces are sort of a weight for the temporal difference error, telling us how much the past states and which feature have contributed  to the learning process.\\
In the evaluation phase of the produced model.pkl, the mean rewards oscillated between -100 and -110.


\paragraph {Resources}
\begin{itemize}
\item Sutton Barto
\item https://github.com/LyWangPX/Reinforcement-Learning-2nd-Edition-by-Sutton-Exercise-Solutions
\end{itemize}

\paragraph {Colleagues I consulted with}

\begin{itemize}
\item Claudio Schiavella
\item Andrea Marzo
\item Nicoló Dentale
\item Charlotte Primiceri
\item Elisa Santini
\end{itemize}
\end{document}